\chapter*{Abstract}

Large language models can learn a task from a handful of labelled examples placed in their prompt, a behaviour called in-context learning.
Which examples are shown matters, and a growing literature selects them by similarity, diversity or learned scores.
Little is known about whether selection helps when the test data shift away from the examples, or when what the examples teach contradicts what the model already believes.

This thesis asks whether choosing the demonstrations can help a Smaller model (7B parameters) follow them over its own priors under controlled distribution shift.
We build a synthetic tabular testbed in which covariate shift, spurious-feature reversal and semantic-prior conflict can each be switched on alone.
Semantic-prior conflict gives the features names whose everyday meaning agrees or disagrees with the rule; unlike concept shift, the rule itself never changes.
On Qwen2.5-7B-Instruct with eight demonstrations, we compare seven hand-designed selection strategies against a zero-shot baseline.
We then train a set-transformer selector (\sata) on ground-truth targets, and look inside the model with the logit lens, linear probes and activation steering.
Every confirmatory test was fixed before its data were seen.

\textbf{Feature names influences the model's predictions more than the choice of demonstrations does.}
Flipping the names costs 0.20 \AUROC (22 of 24 tasks), more than any shift of the data, whereas the hand-designed strategies typically differ from balanced random demonstrations by only 0.025, and none overcame flipped names.
\textbf{A trained selector buys robustness at a cost in predictive performance on the instruction-tuned model}: \sata has no weak environment on a stand-in learner, but on Qwen it evens out the environments by giving up \AUROC in distribution and under covariate shift, and none of its six pre-registered contrasts is significant.
Exploratorily, it helps the base model, whose priors pull less, without that cost.
\textbf{Inside the model, the logit lens and probes suggest that the answer forms only in the last few layers, where the names' effect appears}: the model's own guess is decodable from layer 2 but the demonstrated rule hardly at all, and steering towards the demonstrations moves its outputs without making them more correct.
These tools locate the effect but do not fully explain the decision.
At our scale (7B, eight numeric demonstrations, synthetic tasks), these results bound what selection achieved, not what it can achieve.
